\documentclass[11pt]{article}
\usepackage[T1]{fontenc}
\usepackage[utf8]{inputenc}
\usepackage[a4paper,margin=24mm]{geometry}
\usepackage{amsmath,amssymb}
\usepackage{booktabs}
\usepackage{tabularx}
\usepackage{array}
\usepackage{graphicx}
\usepackage{xcolor}
\usepackage[numbers,sort&compress]{natbib}
\usepackage{microtype}
\usepackage[hidelinks]{hyperref}

\definecolor{maceblue}{HTML}{2F6690}
\definecolor{fnored}{HTML}{C44E52}
\definecolor{lesgreen}{HTML}{4C956C}
\newcommand{\code}[1]{\texttt{#1}}
\newcommand{\macefno}{\textsc{MACE--FNO}}
\newcommand{\meva}{\ensuremath{\mathrm{meV\,\text{\AA}^{-1}}}}
\newcommand{\mevpa}{\ensuremath{\mathrm{meV\,atom^{-1}}}}
\newcommand{\angstrom}{\ensuremath{\text{\AA}}}
\newcommand{\RMSE}{\ensuremath{\mathrm{RMSE}}}
\newcommand{\vect}[1]{\boldsymbol{#1}}
% Values are transcribed from immutable external logs/audits. Keep this file
% small so numerical updates can be reviewed independently from the prose.
\newcommand{\WaterMaceEnergy}{0.986}
\newcommand{\WaterFnoEnergy}{0.556}
\newcommand{\WaterMaceForce}{67.57}
\newcommand{\WaterFnoForce}{52.83}
\newcommand{\WaterEnergyGain}{43.6}
\newcommand{\WaterForceGain}{21.8}
\newcommand{\WaterSpectralExponent}{1.878}
\newcommand{\WaterSpectralFreeRtwo}{0.977}
\newcommand{\WaterSpectralRtwo}{0.973}

\newcommand{\AuMaceEnergy}{2.312}
\newcommand{\AuFnoEnergy}{0.116}
\newcommand{\AuMaceForce}{56.99}
\newcommand{\AuFnoForce}{13.66}

% Replaced after validation-only convergence selection and endpoint rerun.
\newcommand{\AuSelectedUndoped}{pending}
\newcommand{\AuSelectedDoped}{pending}
\newcommand{\AuSelectedSign}{pending}
\newcommand{\AuSelectedConfiguration}{pending}
\newcommand{\AuSelectedObjective}{pending}
\newcommand{\AuSelectionNarrative}{The convergence jobs are in progress.}
\newcommand{\AuSweepRows}{%
Baseline & 24 & 4 & 16 & 4 & 0.06638 & 13.66\\
\multicolumn{7}{c}{Remaining sweep members in progress.}\\}


\title{\textbf{Learning residual interactions with Fourier neural operators: Modeling long-range and nonlocal interactions}}
\author{Wei Chen}
\date{1 September 2026}

\begin{document}
\maketitle

\begin{abstract}
Local equivariant interatomic potentials are accurate and efficient when the
target energy is well approximated by finite-range atomic contributions. They
can nevertheless miss dielectric screening, charge-transfer-mediated
coupling, and other effects whose relevant length scale exceeds their
receptive field. We develop a conservative residual model in which a frozen
MACE potential supplies local invariant descriptors and a Fourier neural
operator (FNO) propagates descriptor-dependent latent fields over a global
mesh. The construction does not assign physical charges: it learns the
energy--force residual left by the local baseline.

On the held-out water--SCAN test set, the current periodic 3D EqGINO-style
model reduces the MACE energy RMSE from \WaterMaceEnergy{} to
\WaterFnoEnergy{}~\mevpa{} and the force RMSE from \WaterMaceForce{} to
\WaterFnoForce{}~\meva{}. A low-wavevector curvature audit finds a fitted
exponent of \WaterSpectralExponent{} and \(R^2=\WaterSpectralRtwo{}\) for a
fixed \(1/k^2\) response, consistent with an electrostatic-like component but
not sufficient to identify the latent fields as charge density. On
Au$_2$/MgO, the established 2D FNO reduces the held-out force RMSE from
\AuMaceForce{} to \AuFnoForce{}~\meva{} and reproduces the DFT reversal of the
wetting preference after remote Al doping. A controlled grid/mode/channel
sweep and a post-selection endpoint rerun are reported in the final version.

We establish feasibility, conservative forces, and sensitivity to a remote
perturbation. We do not yet establish transfer across arbitrary cell shapes,
a unique electrostatic interpretation, Born effective charges, stress, or
production molecular-dynamics efficiency.
\end{abstract}

\section{Motivation}

Message-passing machine-learning interatomic potentials (MLIPs) generally
remain limited in their description of long-range electrostatics, even when
equipped with several interaction layers. Latent Ewald summation
(LES)~\citep{cheng2025les} augments a local equivariant MLIP with learned
latent charge-like variables whose long-range interactions are evaluated
through an explicit Coulomb kernel. qNEP~\citep{fan2026qnep} similarly uses
environment-dependent latent charges to describe long-range electrostatics
and polarization. Here, we ask whether a global operator can learn the
structured error that remains after a strong local potential has already
described short-range chemistry. We use the decomposition
\begin{equation}
 E(\mathcal R,\mathbf H)=E_{\mathrm{MACE}}(\mathcal R,\mathbf H)
 +\Delta E_{\mathrm{FNO}}(\mathcal R,\mathbf H),
 \label{eq:decomposition}
\end{equation}
where \(\mathcal R\) denotes the atomic configuration and \(\mathbf H\) the
cell. We keep MACE frozen during residual training while retaining descriptor
derivatives with respect to coordinates. The FNO correction therefore
contributes conservative forces through the complete computational graph.

We presently support three findings.
\begin{enumerate}
 \item We find that a periodic 3D residual improves both energy and force errors for
 water--SCAN, rather than merely fitting a constant energy offset.
 \item We show that a slab-specific 2D FNO residual matches the remote-Al sign reversal in
 Au$_2$/MgO, a deliberately nonlocal endpoint test that the one-interaction
 local baseline fails.
 \item We find that the 3D learned response is close to \(1/k^2\) over the probed
 long-wavelength modes. This is a useful physical diagnostic; it is not a
 charge-density measurement and is not used as a training loss.
\end{enumerate}

We therefore interpret the current model as combining local chemistry from
frozen MACE with a learned global residual response from the FNO. The method
is broader, but less identifiable, than approaches that predict
latent charges and apply a prescribed Ewald kernel. That distinction is
central to the comparison with LES and qNEP in Sec.~\ref{sec:comparison}.

\section{The nonlocal gap in local MLIPs}

\subsection{Finite receptive fields}

MACE builds equivariant atom-centered representations by message passing within
a radial cutoff and maps them to an extensive energy~\citep{batatia2022mace}.
Higher body order makes the local representation
expressive with few interaction layers, but does not remove the finite
receptive field. Two structures can therefore present nearly identical local
neighborhoods to the model while differing through remote dopants, boundary
conditions, polarization, or collective charge redistribution.

This limitation is not specific to MACE. Long-range-aware atomistic models
have used explicit charges, charge equilibration, Wannier centers, analytical
electrostatics, long-distance descriptors, reciprocal-space messages,
attention, or combinations of these ideas~\citep{grisafi2019longrange,ko2021fourth,zeng2022dplr,
unke2021spookynet,kosmala2023ewaldmp,faller2024density}. These approaches differ mainly in
where they place physical prior information and which latent quantities are
identifiable.

\subsection{Residual learning as a controlled intervention}

We freeze the local model for two purposes. First, this isolates the
incremental information available to a global branch. Second, it prevents an
end-to-end model from silently reallocating short-range chemistry to the FNO,
which would make the role of the global operator harder to diagnose. We do
not assume that the learned correction is purely long ranged. Residual
learning instead makes that hypothesis testable through controlled baselines,
reciprocal-space diagnostics, cell-size tests, and physically targeted
endpoints.

\section{Fourier neural operator construction}

\subsection{Continuous operator and discrete Fourier layer}

A conventional neural network maps one finite-dimensional vector to another.
A neural operator instead maps an input function to an output function,
\(
 \mathcal G_\omega:a(\mathbf r)\mapsto u(\mathbf r)
\), with parameters shared over the spatial domain and, in the continuum
idealization, across discretizations~\citep{kovachki2023neuraloperator}. A
generic operator layer can be written as
\begin{equation}
 (\mathcal K_\ell v_\ell)(\mathbf r)
 =\int_\Omega
 \kappa_\ell(\mathbf r,\mathbf r')v_\ell(\mathbf r')\,\mathrm d\mathbf r'.
 \label{eq:integral-operator}
\end{equation}
When the domain is periodic and the kernel is translation invariant,
\(\kappa_\ell(\mathbf r,\mathbf r')=\kappa_\ell(\mathbf r-\mathbf r')\), the
convolution theorem converts Eq.~\eqref{eq:integral-operator} into a
multiplication in reciprocal space. The Fourier neural operator (FNO)
parameterizes that multiplier rather than storing a dense real-space kernel
between every pair of mesh points~\citep{li2021fno}.

On a mesh with multi-index \(\mathbf m=(m_1,\ldots,m_d)\) and dimensions
\(N_1\times\cdots\times N_d\), we use the discrete transform
\begin{equation}
 \widehat v_{\ell a}(\mathbf n)
 =\sum_{\mathbf m}v_{\ell a}(\mathbf m)
 \exp\!\left[-2\pi i\sum_{\alpha=1}^{d}
 \frac{n_\alpha m_\alpha}{N_\alpha}\right],
 \label{eq:discrete-fourier}
\end{equation}
where \(a\) labels a feature channel and \(\mathbf n\) labels an integer
reciprocal-grid mode. One nonlinear FNO block is
\begin{equation}
 v_{\ell+1}(\mathbf r)=\sigma\!\left[
 W_\ell v_\ell(\mathbf r)+
 \mathcal F^{-1}\{R_\ell(\mathbf n)\widehat v_\ell(\mathbf n)\}(\mathbf r)
 \right].
 \label{eq:fno-layer}
\end{equation}
The matrix \(W_\ell\) mixes channels independently at every mesh point. The
spectral matrix \(R_\ell(\mathbf n)\) mixes input and output channels at each
retained mode, and all unretained spectral coefficients are set to zero. A
single retained basis function spans the entire periodic domain. The FNO
therefore obtains a global receptive field in one block; increasing the mode
cutoff adds shorter-wavelength resolution rather than increasing the maximum
communication distance.

Our nonlinear operators first lift the \(C\) deposited source channels to
\(H_f\) hidden channels, apply \(L_f\) blocks of the form in
Eq.~\eqref{eq:fno-layer}, and project back to \(C\) response channels:
\begin{equation}
 v_0=P\boldsymbol\rho,\qquad
 v_{\ell+1}=\mathcal B_\ell(v_\ell),\qquad
 \boldsymbol\phi=Q_2\,\mathrm{GELU}(Q_1v_{L_f}).
 \label{eq:fno-stack}
\end{equation}
The lifting and projection maps are pointwise \(1\times1\) convolutions. We
disable their biases and do not add Cartesian position embeddings. These
choices preserve discrete periodic translation equivariance and prevent the
residual energy from acquiring a configuration-independent offset through
the mesh branch. The local pathway in Eq.~\eqref{eq:fno-layer} also preserves
mesh information outside the retained spectral band, although only the
spectral pathway communicates it globally.

Modes and channels control different kinds of capacity. Modes specify which
spatial wavelengths receive learned global mixing. Channels specify how many
components of a field are carried at one mesh point. We distinguish the
\(C\) latent source channels \(\rho_c\) from the \(H_f\) hidden channels
inside the operator. Increasing \(C\) gives the model more independent
source--response components, whereas increasing \(H_f\) widens every FNO
block. Neither index labels a chemical species, atomic orbital, charge state,
or electrostatic multipole.

\subsection{Invariant atomic sources}

The FNO does not receive atomic numbers or coordinates as isolated scalar
inputs. We first run the frozen MACE model and collect all even-parity scalar
\(0e\) components from the product output of each interaction layer. Their
concatenation defines the invariant descriptor \(\mathbf h_i\) for atom
\(i\). This choice retains local chemical information while ensuring that the
learned source amplitudes are unchanged by a rigid rotation of the atomic
configuration.

We map each descriptor to \(C\) raw scalar amplitudes with
\begin{equation}
 \widetilde{\mathbf s}_i
 =A_2\,\mathrm{SiLU}\!\left[A_1\,\mathrm{LN}(\mathbf h_i)\right],\qquad
 \mathbf s_i=\widetilde{\mathbf s}_i
 -\frac{1}{N_g}\sum_{j\in g}\widetilde{\mathbf s}_j .
 \label{eq:neutral-source}
\end{equation}
Here \(A_1\) expands to the source-head hidden width, \(A_2\) projects to
\(C\) channels without a bias, and the second operation is performed
separately for every structure \(g\) and channel. Consequently,
\(\sum_{i\in g}s_{ic}=0\) to numerical precision. Neutralization removes the
unconstrained zero-wavevector source and is appropriate for the neutral
systems considered here. It is not a treatment of a physically charged cell
or of the compensating background used by periodic electronic-structure
codes.

Equation~\eqref{eq:neutral-source} also clarifies how the FNO uses MACE. MACE
determines which local environments generate the latent sources; the mesh
operator determines how those sources interact across the full cell. Remote
information does not enter \(\mathbf h_i\) directly when it lies outside the
MACE receptive field. It reaches atom \(i\) through the response field
generated by sources deposited elsewhere.

\subsection{Cubic B-spline particle--mesh assignment}

We express each atomic position in fractional cell coordinates and convert a
fractional component to a mesh coordinate \(u=p+t\), where \(p=\lfloor
u\rfloor\) and \(0\leq t<1\). The four cardinal cubic B-spline weights are
\begin{equation}
 \mathbf b(t)=\frac{1}{6}\left(
 (1-t)^3,
 3t^3-6t^2+4,
 -3t^3+3t^2+3t+1,
 t^3\right),
 \label{eq:bspline-weights}
\end{equation}
and correspond to mesh offsets \((-1,0,1,2)\). Their tensor product assigns
one atom to 16 points on a planar mesh or 64 points on a three-dimensional
mesh. This compact support is smooth through second derivatives, which is
important because force training differentiates through the assignment.

For structure \(g\), the deposited field is
\begin{equation}
 \rho_{gc}(\mathbf m)
 =\frac{1}{\Delta\Omega_g}
 \sum_{i\in g}s_{ic}\,A_i(\mathbf m),\qquad
 A_i(\mathbf m)=\prod_{\alpha\in\mathcal P}
 b_{m_\alpha-p_{i\alpha}}(t_{i\alpha}),
 \label{eq:latent-source-field}
\end{equation}
where \(\mathcal P\) contains the meshed coordinates and a spline factor is
zero outside its four-point support. The element \(\Delta\Omega_g\) is the
mesh-point area for a planar field and the voxel volume for a slab or bulk
field. This normalization makes the assignment conservative:
\begin{equation}
 \Delta\Omega_g\sum_{\mathbf m}\rho_{gc}(\mathbf m)
 =\sum_{i\in g}s_{ic}=0.
 \label{eq:mesh-neutrality}
\end{equation}
The construction follows the smooth particle--mesh assignment used in
particle--mesh electrostatics~\citep{darden1993pme,essmann1995pme}, but the
deposited amplitudes here are learned latent variables.

The symbol \(\rho_c\) therefore denotes a real-space latent source field, not
an electronic, induced, or nuclear charge density. The channel basis can be
rotated or rescaled if compensating changes are made in the source head and
field operator. Energy--force training alone does not fix physical charge
units, so moments of \(\rho_c\) cannot presently be interpreted as
polarization or Born effective charges.

\subsection{Response field, scalar energy and conservative forces}

The FNO maps the deposited sources to response fields,
\(\boldsymbol\phi=\mathcal G_\omega[\boldsymbol\rho;\mathbf H]\). We then
construct one scalar correction for each structure:
\begin{equation}
 \Delta E_{g,\mathrm{FNO}}
 =\frac{\Delta\Omega_g}{2}
 \sum_{c,\mathbf m}\rho_{gc}(\mathbf m)\phi_{gc}(\mathbf m)
 \longrightarrow
 \frac{1}{2}\sum_c\int_{\Omega_g}
 \rho_c(\mathbf r)\phi_c(\mathbf r)\,\mathrm d\mathbf r.
 \label{eq:mesh-energy}
\end{equation}
This form is bilinear in the displayed source and response fields, but it is
not generally quadratic in \(\boldsymbol\rho\) because the nonlinear FNO
makes \(\boldsymbol\phi\) depend nonlinearly on
\(\boldsymbol\rho\). It gives an extensive scalar with the correct mesh
measure for each cell and provides a single energy from which all residual
forces are derived.

The term ``response field'' is deliberate. For a nonlinear or nonsymmetric
operator, \(\boldsymbol\phi\) is not necessarily the functional derivative of
the energy. In compact notation,
\begin{equation}
 \frac{\delta\Delta E_{\mathrm{FNO}}}{\delta\boldsymbol\rho}
 =\frac{1}{2}\left[
 \boldsymbol\phi+
 J_{\mathcal G}(\boldsymbol\rho)^{\mathsf T}\boldsymbol\rho
 \right],
 \label{eq:energy-functional-derivative}
\end{equation}
where \(J_{\mathcal G}\) is the Jacobian of the learned field operator. Only a
linear self-adjoint operator reduces Eq.~\eqref{eq:energy-functional-derivative}
to the usual electrostatic potential. Our current architecture does not impose
linearity, self-adjointness, or Poisson's equation.

The total force follows from the summed MACE and mesh energies,
\begin{equation}
 \mathbf F_i=-\frac{\partial}{\partial\mathbf r_i}
 \left(E_{\mathrm{MACE}}+\Delta E_{\mathrm{FNO}}\right).
 \label{eq:total-force}
\end{equation}
The coordinate derivative of the mesh field contains a direct assignment path
and a descriptor-mediated path,
\begin{equation}
 \frac{\partial\boldsymbol\rho}{\partial\mathbf r_i}
 =\left.\frac{\partial\boldsymbol\rho}{\partial\mathbf r_i}\right|_{\mathbf s}
 +\sum_{j,q}
 \frac{\partial\boldsymbol\rho}{\partial\mathbf s_j}
 \frac{\partial\mathbf s_j}{\partial\widetilde{\mathbf s}_q}
 \frac{\partial\widetilde{\mathbf s}_q}{\partial\mathbf h_q}
 \frac{\partial\mathbf h_q}{\partial\mathbf r_i}.
 \label{eq:force-paths}
\end{equation}
The first term moves the spline stencil. The second changes the latent source
amplitudes as local MACE environments change; the graph-wise mean in
Eq.~\eqref{eq:neutral-source} couples the source derivatives within each
structure. We freeze all MACE parameters but deliberately retain the
coordinate graph through \(\mathbf h_i\). Freezing the weights therefore
does not detach the descriptors or remove their contribution to the residual
forces.

The FNO consequently learns a configuration-dependent source--response map,
not a direct scalar regressor and not a single fixed reciprocal-space kernel.
Its local dependence enters through \(g_\psi(\mathbf h_i)\); its nonlocal
dependence enters through the global Fourier blocks. The low-wavevector audit
later in this report probes the curvature of Eq.~\eqref{eq:mesh-energy} under
controlled Fourier perturbations. A \(1/k^2\) curvature is evidence for an
electrostatic-like component of the learned energy, but it does not identify
any one latent channel as charge.

Figure~\ref{fig:fno-workflow} summarizes the complete forward and derivative
paths. It distinguishes neutral latent fields from physical charge density and
fixed MACE weights from the coordinate derivatives that remain active.

\begin{figure}[t]
\centering
\includegraphics[width=\textwidth]{fno_workflow.pdf}
\caption{Information flow in our frozen-MACE residual architecture. MACE
provides the local baseline energy and invariant descriptors while its weights
remain fixed. The learned branch deposits neutral latent sources, propagates
them with a geometry-aware FNO, and returns a scalar residual energy. Automatic
differentiation of the summed energy gives conservative forces. The latent
mesh fields are learned internal variables, not assigned charge densities.}
\label{fig:fno-workflow}
\end{figure}

\section{Implemented geometries, symmetries and cost}

\subsection{Slab-resolved 2D FNO}

For a supported cluster, the physical boundary conditions are periodic in the
surface plane and finite through the vacuum direction. A fully periodic 3D
FFT would spuriously identify the top and bottom of the finite slab window. We
define the surface plane with the first two cell vectors, wrap the associated
fractional coordinates, and project Cartesian positions onto their unit
normal. The finite normal coordinate is not wrapped. We center a window of
length \(L_z\) on the instantaneous mean atomic height, which makes a rigid
translation normal to the surface exactly invariant. Cubic-spline support
that crosses a window boundary is accumulated into the nearest boundary
voxel rather than wrapped to the opposite surface.

We retain a field of shape \((C,n_z,n_x,n_y)\), transform only \(x,y\), and
mix \(z\) explicitly. Each nonlinear block is
\begin{equation}
 v_{\ell+1}=\mathrm{GELU}\left[
 \mathcal F_{xy}^{-1}R_\ell(\mathbf k_\parallel)\mathcal F_{xy}v_\ell+
 Z_\ell v_\ell+W_\ell v_\ell\right].
 \label{eq:slab-fno-layer}
\end{equation}
The first term applies the same learned 2D spectral convolution independently
to each normal layer. In the established Au$_2$/MgO model, \(Z_\ell\) is an
independent dense \(n_z\times n_z\) matrix for every hidden channel, shared
over the lateral mesh. It couples any two normal layers without identifying
the two boundaries. The package also provides zero-padded local normal
convolutions and a controlled linear operator with an explicit dense
\(R(\mathbf k_\parallel,z,z')\), but neither is used for the reported endpoint
result.

We call this architecture the 2D FNO because its Fourier propagation spans the
two periodic in-plane directions. The normal coordinate remains explicitly
resolved on a finite, nonperiodic grid. A purely planar projection with field
shape \((C,n_x,n_y)\) is retained in the code as an ablation, but the reported
Au$_2$/MgO model uses the slab-resolved field in Eq.~\eqref{eq:slab-fno-layer}.

The electrostatic reference is different from the 3D-periodic \(1/k^2\)
kernel. After Fourier transforming only the surface plane, an isolated-slab
Coulomb Green's function has the form
\begin{equation}
 G_{2\mathrm D}(k_\parallel;z,z')\propto
 \frac{\exp[-k_\parallel|z-z'|]}{k_\parallel},\qquad k_\parallel>0.
 \label{eq:slab-green-function}
\end{equation}
The 2D FNO spectral diagnostic therefore tests a \(1/k_\parallel\) planar
monopole response and the associated finite-\(z\) profile, rather than applying
the bulk \(1/k^2\) criterion.

The spline assignment is exact under an integer mesh translation but only
approximately invariant under an arbitrary continuous translation. With
\code{lateral-interlacing=2}, we average energies from the four combinations
of zero and half-grid shifts along the two periodic directions:
\begin{equation}
 \Delta E_{\mathrm{int}}
 =\frac{1}{4}\sum_{\delta_x,\delta_y\in\{0,1/2\}}
 \Delta E(\delta_x,\delta_y).
 \label{eq:slab-interlacing}
\end{equation}
The four replicas are evaluated as one larger batch. Interlacing reduces the
particle--mesh egg-box error, but does not eliminate it and does not yield one
unique deposited field for spectral analysis.

For a square lateral grid, we implement the four rotations and four
reflections of \(D_4\). During training, successive forward passes cycle
through one transformed image, providing balanced symmetry augmentation
without an eightfold cost. During evaluation, we use the group average
\begin{equation}
 \overline{\mathcal G}(\rho)
 =\frac{1}{8}\sum_{g\in D_4}
 T_g^{-1}\mathcal G(T_g\rho),
 \label{eq:d4-average}
\end{equation}
which makes the evaluated field operator exactly \(D_4\)-equivariant up to
floating-point precision. The established Au$_2$/MgO model combines this
evaluation-time average with four-origin lateral interlacing.

\subsection{Periodic 3D and the EqGINO-style contraction}

For bulk systems, we convert Cartesian positions with the full cell matrix,
wrap all three fractional coordinates, and deposit each source on 64 periodic
mesh points. The voxel measure in Eq.~\eqref{eq:mesh-energy} is
\(|\det\mathbf H|/(n_zn_xn_y)\). The conventional 3D FNO uses a real-to-complex
FFT and learns independent complex channel-mixing matrices for the retained
positive and negative mode sectors. This implementation supports a fixed
general cell, but it does not by itself provide transfer between arbitrary
cell shapes because the learned weights are indexed by fractional-grid modes.

For cubic bulk water we use a more restrictive EqGINO-style contraction. We
share one real channel-mixing matrix among integer reciprocal vectors with the
same squared radius,
\begin{equation}
 R_\ell(\mathbf n)=R_{\ell,s},\qquad
 s=n_x^2+n_y^2+n_z^2.
 \label{eq:eqgino-shell}
\end{equation}
All signed permutations of the cubic axes leave \(s\) unchanged. The spectral
convolution is therefore exactly equivariant under the 48 signed axis
permutations of the cubic grid. Real shell weights preserve Hermitian symmetry
for the real scalar fields used here. This is an atomistic real-to-real
adaptation of EqGINO~\citep{kim2026eqgino}, rather than a reproduction of its
complex-feature architecture.

Equation~\eqref{eq:eqgino-shell} replaces one matrix per retained wavevector by
one matrix per distinct radial shell. We additionally divide the hidden
channels into \(G\) spectral groups and apply a block-diagonal contraction.
The pointwise \(W_\ell\) pathway remains dense and communicates between groups
at each mesh point. For the water configuration, four groups and radial-shell
sharing reduce the residual parameter count from 271,824 for the matched
dense-wavevector FNO to 12,112, while preserving two nonlinear Fourier blocks.

EqGINO shell sharing requires a cubic mesh, equal mode cutoffs, and cubic
cells. The water cells are aligned cubes with varying lengths. We append
\begin{equation}
 \chi_g=\log L_g,\qquad
 L_g=|\det\mathbf H_g|^{1/3},
 \label{eq:cell-conditioning}
\end{equation}
as a spatially constant input channel so that the nonlinear operator receives
an explicit physical length scale. The current cell validator accepts only
positive uniform scalings of the reference cube in this mode. Arbitrary
anisotropic cell changes require metric-aware wavevectors or separate
shape-compatible models.

Three-dimensional volume interlacing is also implemented. At level two, it
averages the eight combinations of zero and half-grid offsets in the three
periodic directions. We disable it in the reported water fit for throughput,
which leaves the continuous-translation error quantified in the numerical
audit.

\subsection{Computational scaling and present runtime tradeoffs}

Let \(N\) be the number of atoms, \(M=n_zn_xn_y\) the number of mesh points,
\(H_f\) the hidden width, \(K\) the number of retained modes, \(L_f\) the
number of FNO blocks, and \(G\) the number of spectral groups. Cubic B-spline
assignment costs \(O(64NC)\) in 3D and \(O(64NC)\) for the slab-resolved mesh;
the purely planar projection costs \(O(16NC)\). The leading operator costs are
approximately
\begin{equation}
 O(L_fH_fM\log M)
 +O(L_fKH_f^2/G)
 +O(L_fMH_f^2),
 \label{eq:fno-cost}
\end{equation}
from the FFTs, grouped spectral contractions, and dense pointwise channel
mixing, respectively. The slab model additionally incurs
\(O(L_fH_fn_z^2n_xn_y)\) for global normal mixing. Memory scales primarily as
\(O(H_fM)\) per structure, apart from temporary symmetry or interlacing
replicas.

At fixed mesh size, the FNO communication range is global and its cost does
not depend on the separation between two atoms. For increasing physical cell
size at fixed spatial resolution, however, \(M\) grows with volume and the FFT
term approaches \(O(M\log M)\). The method is therefore more favorable than a
dense all-pairs interaction, but less local and less communication-friendly
than the frozen MACE branch.

Symmetry restoration and interlacing introduce transparent multipliers. The
water calculation uses neither volume interlacing nor explicit cubic group
averaging because EqGINO is equivariant by construction. The Au$_2$/MgO
inference path evaluates four mesh origins and eight \(D_4\) images, equivalent
to 32 FNO replicas even though they are vectorized into large batches. During
training, cycling one \(D_4\) image reduces this to four interlaced replicas.
These choices improve symmetry and translation behavior but explain why the
slab-resolved model is slower than an unaveraged 2D FNO.

Freezing MACE reduces trainable parameters and permits a one-interaction local
baseline, but it does not make the backbone forward or coordinate derivatives
free. Force training also differentiates a force loss through
Eq.~\eqref{eq:total-force}, producing mixed derivatives with respect to model
parameters and atomic coordinates. We therefore do not yet claim that
MACE--FNO inference is faster than a deeper standalone MACE model. A fair
runtime comparison must report energy-only and energy--force throughput,
batch size, mesh resolution, symmetry averaging, interlacing, precision, and
system size.

\clearpage
\subsection{Exact symmetries and numerical approximations}

The current implementation separates architectural guarantees from numerical
particle--mesh approximations. Table~\ref{tab:symmetry-status} summarizes this
distinction.

\begin{table}[h!]
\centering
\caption{Symmetries and conservation properties of the implemented residual
models. ``Exact'' denotes equality up to floating-point precision for the
stated geometry.}
\label{tab:symmetry-status}
\small
\begin{tabularx}{\textwidth}{>{\raggedright\arraybackslash}p{31mm}
>{\raggedright\arraybackslash}X>{\raggedright\arraybackslash}p{29mm}}
\toprule
Property & Implementation mechanism & Status\\
\midrule
Latent neutrality & Graph-wise channel mean subtraction & Exact\\
Energy conservation & One scalar mesh functional differentiated for forces & Exact\\
Periodic lattice shift & Fractional wrapping and periodic spline indices & Exact\\
Integer mesh shift & Translation-equivariant FFT and pointwise paths & Exact\\
Rigid slab-normal shift & Mean-height finite-window center & Exact\\
Slab \(D_4\) operations & Evaluation-time group average & Exact at evaluation\\
Cubic signed permutations & EqGINO radial-shell sharing & Exact for cubic cells\\
Arbitrary translation & Finite particle--mesh origin & Approximate; audited\\
Arbitrary cell-shape transfer & No general metric conditioning & Not supported\\
\bottomrule
\end{tabularx}
\end{table}

We test these statements through source sums, force additivity, the acoustic
sum rule, finite differences, full-lattice and one-grid translations, and the
promised point-group operations. Finite-difference agreement establishes that
the forces are derivatives of the discretized energy. It does not establish
continuous translation invariance, because atom-to-mesh assignment still
depends on the grid origin. Interlacing suppresses this egg-box error at an
explicit computational cost.

\section{Training and model selection}

We precompute the residual labels of the frozen baseline,
\begin{equation}
 \Delta E_g^{\mathrm{ref}}=E_g^{\mathrm{ref}}-E_g^{\mathrm{MACE}},
 \qquad
 \Delta\mathbf F_i^{\mathrm{ref}}
 =\mathbf F_i^{\mathrm{ref}}-\mathbf F_i^{\mathrm{MACE}}.
 \label{eq:residual-targets}
\end{equation}
Caching these labels avoids recomputing frozen-MACE reference forces for every
optimization batch. We still execute the MACE forward graph because the
current atomic descriptors and their coordinate derivatives are required by
Eq.~\eqref{eq:force-paths}.

We optimize the source head and FNO parameters against residual energies and
forces. For a batch of \(B\) structures containing \(N_{\mathrm{tot}}\) atoms,
the training loss is
\begin{equation}
 \mathcal L_{\mathrm{train}}=
 \frac{w_E}{B}\sum_{g=1}^{B}
 \left[
 \frac{\Delta E_g-\Delta E_g^{\mathrm{ref}}}{N_g\sigma_E}
 \right]^2
 +\frac{w_F}{3N_{\mathrm{tot}}}
 \sum_{i,\alpha}
 \left[
 \frac{\Delta F_{i\alpha}-\Delta F_{i\alpha}^{\mathrm{ref}}}{\sigma_F}
 \right]^2.
 \label{eq:training-loss}
\end{equation}
The energy term is an MSE of per-atom residuals, while the force term is an
MSE over Cartesian components. Backpropagating the second term differentiates
the energy gradient in Eq.~\eqref{eq:total-force} with respect to residual
parameters. This mixed derivative is substantially more expensive than
energy-only fitting and is one reason that force convergence controls the
training time.

Both selected benchmarks use \(w_E=w_F=1\). For water--SCAN, we set
\(\sigma_E=0.01\)~eV atom\(^{-1}\) and
\(\sigma_F=0.10\)~eV \angstrom\(^{-1}\). For Au$_2$/MgO, we use the frozen
baseline training errors,
\(\sigma_E=0.002260\)~eV atom\(^{-1}\) and
\(\sigma_F=0.056280\)~eV \angstrom\(^{-1}\). The scales make the two loss
terms dimensionless and prevent their numerical units from setting the
relative importance of energies and forces.

Checkpoint selection applies the same normalization to validation RMSEs:
\begin{equation}
 \mathcal L_{\mathrm{val}}=
 w_E\left(\frac{\RMSE_E}{\sigma_E}\right)^2+
 w_F\left(\frac{\RMSE_F}{\sigma_F}\right)^2 .
 \label{eq:validation-objective}
\end{equation}
We never use the reported test sets or Au$_2$/MgO wetting endpoints to choose
a checkpoint or a convergence-sweep configuration.

For the Au$_2$/MgO sweep, we fixed a selection rule before inspecting
endpoints: minimize Eq.~\eqref{eq:validation-objective}; regard configurations
within 1\% of the best value as tied; among tied configurations select the
lower-cost model. We order cost first by the number of trainable residual
parameters and then, if still tied, by the number of mesh points. We vary one
factor at a time around
\((n_x,n_y,n_z)=(24,24,16)\), an in-plane mode cutoff of four, and four latent
source channels. We hold seed 17, 20\,000 optimization steps, float64,
global \(z\) mixing, mean-\(z\) centering, two-way lateral interlacing and \(D_4\)
averaging fixed. We treat this as a convergence study, not a statistical
estimate over random seeds.

\section{Benchmark I: periodic water--SCAN}
\label{sec:water}

\subsection{Protocol}

We use the 128-water-molecule SCAN~\citep{sun2015scan} configurations employed
in the qNEP study~\citep{fan2026qnep}. Each structure contains 384 atoms in an
aligned cubic cell. The published split supplies 1\,388 training and 500 test
structures. We deterministically reserve 139 of the published training
structures for validation, leaving 1\,249 for fitting; the official
500-structure test is held out.

We freeze a one-interaction MACE baseline with \(r_{\max}=4.5\)~\angstrom.
Our residual uses a \(24^3\) mesh, a mode cutoff of four per direction, four
latent source channels, and two FNO blocks with 16 hidden channels. We use
EqGINO shell sharing with four spectral groups, float64 arithmetic, and no
volume interlacing.
We train for 60\,000 steps with validation every 500 steps and select the
checkpoint at step 57\,500. Our residual branch has 12\,112 trainable
parameters; a matched dense-wavevector 3D FNO has 271\,824, so the
EqGINO-style contraction reduces this count by a factor of 22.4 for this
configuration.
We measured 5\,377~s for the complete 60\,000-step job on one NVIDIA A100 GPU.


\subsection{Accuracy}

\begin{table}[t]
\centering
\caption{Held-out water--SCAN errors. Our MACE and MACE+FNO rows use the same
500-structure test set. Published NEP/qNEP values are reproduced from
Fig.~2 of Ref.~\citep{fan2026qnep}; they use a different baseline architecture
and training protocol and are contextual, not a controlled head-to-head
ranking.}
\label{tab:water-results}
\small
\begin{tabularx}{\textwidth}{>{\raggedright\arraybackslash}Xrrr}
\toprule
Model & Energy RMSE (\mevpa) & Force RMSE (\meva) & Stress RMSE (MPa)\\
\midrule
Frozen one-layer MACE (this work) & \WaterMaceEnergy & \WaterMaceForce & --\\
MACE+3D FNO (this work) & \WaterFnoEnergy & \WaterFnoForce & --\\
\midrule
NEP (published qNEP study) & 1.1 & 81 & 219\\
qNEP mode 1 (published) & 0.9 & 73 & 179\\
qNEP mode 2 (published) & 0.8 & 70 & 204\\
\bottomrule
\end{tabularx}
\end{table}

Our residual lowers the held-out energy RMSE by \WaterEnergyGain\% and the
force RMSE by \WaterForceGain\%. A fitted global energy offset changes the
baseline energy RMSE only from 0.986 to 0.971~\mevpa, so the observed
reduction to 0.556~\mevpa{} is not explained by a constant bias. The force
reduction is particularly important because forces are derivatives of the
learned correction and therefore test the full descriptor--mesh pathway.

The qNEP comparison answers a narrower question. In the published study,
adding dynamic-charge electrostatics to NEP reduces energy RMSE by about 27\%
and force RMSE by about 14\% when comparing NEP with qNEP mode 2. Our
relative reductions are larger, and our final errors are smaller, but this
does not establish that FNO is intrinsically more accurate: our local model,
validation split, optimizer, and parameterization differ. A defensible
comparison must also note that the qNEP study supplied Born-effective-charge
targets for 194 water structures to learn dielectric response, whereas our
FNO fit uses only energies and forces and does not predict electrical
observables. A defensible conclusion is that both studies find measurable energy and force information
beyond their corresponding local baselines on this data set.

\subsection{Spectral and numerical audit}

We probed the learned energy curvature by adding neutral Fourier perturbations
to the post-deposition latent field. Over the three lowest
physical \(k\) shells, the dominant positive response follows \(k^{-p}\) with
\(p=\WaterSpectralExponent\) and free-fit \(R^2=\WaterSpectralFreeRtwo\).
Holding \(p=2\) gives \(R^2=\WaterSpectralRtwo\). Over all modes through
integer index two, the free exponent is 1.989 and the fixed-\(1/k^2\)
log-space \(R^2\) is 0.981. An anisotropic fit
\(1/(\mathbf k^\mathsf T B\mathbf k)\) yields a trace-normalized \(B\)
within \(1.2\times10^{-4}\) of the identity on its diagonal and below
\(1.1\times10^{-4}\) off diagonal, as expected for cubic shell sharing.

\begin{table}[t]
\centering
\caption{Selected numerical checks for the water 3D model. Error metrics in
the first two rows use a fixed 32-structure audit subset; complete held-out
metrics are in Table~\ref{tab:water-results}.}
\label{tab:water-audit}
\begin{tabularx}{\textwidth}{>{\raggedright\arraybackslash}Xrr}
\toprule
Quantity & Frozen MACE & MACE+FNO\\
\midrule
Subset energy RMSE (\mevpa) & 0.6248 & 0.2856\\
Subset force RMSE (\meva) & 53.8488 & 43.7448\\
\midrule
\multicolumn{2}{l}{Maximum latent source sum} & \(4.81\times10^{-13}\)\\
\multicolumn{2}{l}{Finite-difference force discrepancy (\meva)} &
\(5.90\times10^{-6}\)\\
\multicolumn{2}{l}{One-grid translation energy change (meV)} &
\(1.00\times10^{-11}\)\\
\multicolumn{2}{l}{Largest discrete cubic energy change (meV)} &
\(1.82\times10^{-11}\)\\
\multicolumn{2}{l}{Largest 0.1-\angstrom{} arbitrary translation change (meV)} &
20.71\\
\bottomrule
\end{tabularx}
\end{table}

Our exact promised checks pass at numerical precision, but an arbitrary
0.1-\angstrom{} rigid translation can change the residual energy by up to
20.7~meV. This is the clearest current numerical limitation of the
non-interlaced 3D model. Finite-difference agreement shows that the reported
forces remain the correct derivatives of this discretized energy; it does
not remove the grid-origin error.

The qNEP paper reports that NEP and qNEP radial distribution functions are
essentially indistinguishable for liquid water. Accordingly, we do not claim
that an FNO correction should produce a visibly different pair distribution.
We instead identify dipole correlations, dielectric response, infrared
spectra, finite-size scaling, and field response as more discriminating
dynamical targets, although our present latent fields cannot yet predict
these as physical observables.

\section{\texorpdfstring{Benchmark II: Au$_2$/MgO with remote Al dopants}{Benchmark II: Au2/MgO with remote Al dopants}}
\label{sec:aumgo}

\subsection{A discriminating test of nonlocal response}

The Au$_2$/MgO benchmark contains undoped slabs and slabs in which three Al
substitutions lie in the fifth subsurface layer, more than 10~\angstrom{} from
the adsorbed Au dimer~\citep{king2025charges,kim2025universal}. Local Au environments can
remain nearly unchanged while the remote dopants reverse the relative
stability of non-wetting and wetting endpoints. With
\begin{equation}
 \Delta E_{\mathrm{wet}}=
 E(\mathrm{wetting})-E(\mathrm{non\mbox{-}wetting}),
\end{equation}
DFT gives a positive value for undoped MgO and a negative value after doping.
This sign reversal is a more relevant test of nonlocal sensitivity than a
single aggregate RMSE.

We use 4\,500 fit/validation structures and 500 independent test structures,
and we fix a 225-structure validation subset within the former. We freeze a
one-interaction MACE model with \(r_{\max}=5.5\)~\angstrom. Our established
2D FNO residual uses a
\(24\times24\times16\) mesh, an in-plane mode cutoff of four, and four latent
source channels. It contains two FNO blocks with 16 hidden channels, global
finite-\(z\) mixing, two-way lateral interlacing, \(D_4\) averaging and float64
arithmetic, with 50\,624 trainable
parameters. We measured 1~h 13~min for the corresponding 20\,000-step job on
one A100 GPU.


\subsection{Established test accuracy and endpoint result}

\begin{table}[t]
\centering
\caption{Held-out Au$_2$/MgO errors before the present convergence sweep.}
\label{tab:aumgo-errors}
\begin{tabular}{lrr}
\toprule
Model & Energy RMSE (\mevpa) & Force RMSE (\meva)\\
\midrule
Frozen one-interaction MACE & \AuMaceEnergy & \AuMaceForce\\
MACE+2D FNO & \AuFnoEnergy & \AuFnoForce\\
\bottomrule
\end{tabular}
\end{table}

Our residual reduces energy RMSE by about 95\% and force RMSE by about 76\%.
Our 32-structure strict audit gives an energy RMSE of 0.108~\mevpa{} and
force RMSE of 12.86~\meva. Source neutrality, force finite differences,
one-grid lateral translations, and \(D_4\) transformations pass. The remaining
0.1-\angstrom{} lateral translation error is 1.36~meV in energy, with a
3.02~\meva{} RMS net residual-force component along the less favorable
lateral direction. Interlacing substantially reduces this egg-box effect
relative to an unaveraged mesh but does not make continuous translations
exact.

\begin{table}[t]
\centering
\caption{Remote-Al wetting preference in meV. MACE and MACE--LES values are
published in Ref.~\citep{kim2025universal}; the established FNO row is our
independent endpoint relaxation using the same definition. Match asks whether
the undoped-to-doped sign reversal is recovered.}
\label{tab:aumgo-switch}
\begin{tabular}{lrrc}
\toprule
Model & Undoped & Al-doped & DFT sign reversal\\
\midrule
DFT & 934.8 & -66.9 & reference\\
One-interaction MACE & 419.0 & 419.0 & no\\
MACE--LES & 919.3 & -56.5 & yes\\
MACE+2D FNO (established) & 924.5 & -55.4 & yes\\
MACE+2D FNO (selected sweep) & \AuSelectedUndoped & \AuSelectedDoped &
\AuSelectedSign\\
\bottomrule
\end{tabular}
\end{table}

Our established FNO is within 10.4~meV of DFT for the undoped splitting and
11.5~meV for the doped splitting, and is also close to the published MACE--LES
results. We take this as strong evidence that the global residual branch is
sensitive to the remote substitution. We do not yet claim that FNO
outperforms LES: the benchmark contains only one dopant pattern and two
relaxed endpoint differences, and the FNO model is computationally richer.

\subsection{Grid, Fourier-mode and source-channel convergence}

Here \(n_{\mathrm{mode}}\) is the in-plane Fourier cutoff and \(C\) is the
number of latent source fields. The hidden width remains fixed at 16 channels.

\begin{table}[t]
\centering
\caption{Validation-selected convergence of the 2D FNO grid, Fourier-mode
cutoff and latent source-channel count. All members differ by one factor from
the established configuration. Test metrics are revealed only after applying
the pre-registered validation rule.}
\label{tab:aumgo-convergence}
\begin{tabular}{lrrrrrr}
\toprule
Configuration & \(n_{xy}\) & \(n_{\mathrm{mode}}\) & \(n_z\) & \(C\) &
\(\mathcal L_{\mathrm{val}}\) & Test F RMSE\\
\midrule
\AuSweepRows
\bottomrule
\end{tabular}
\end{table}

We select \AuSelectedConfiguration, with validation objective
\AuSelectedObjective. \AuSelectionNarrative

\section{Relation to qNEP and LES}
\label{sec:comparison}

\subsection{Different priors, not merely different software}

LES starts from a strong prior: a low-dimensional latent source learned from
local features interacts through a prescribed Ewald
kernel~\citep{cheng2025les}. In a simplified periodic form,
\begin{equation}
 q_i=Q_\phi(\mathbf h_i),\quad
 S(\mathbf k)=\sum_iq_i e^{i\mathbf k\cdot\mathbf r_i},\quad
 E_{\mathrm{LR}}\propto
 \sum_{\mathbf k\ne0}\frac{e^{-\sigma^2k^2/2}}{k^2}|S(\mathbf k)|^2.
\end{equation}
The universal LES implementation augments several local MLIPs and can infer
polarization and Born effective charges from energy--force
training~\citep{king2025charges,kim2025universal,zhong2025electrical}.

qNEP also predicts environment-dependent charges but couples them to explicit
electrostatics implemented with Ewald summation or particle--particle
particle--mesh methods~\citep{fan2026qnep}. Mode~1 includes the real-space,
reciprocal-space and self-energy Ewald terms, whereas mode~2 retains only the
reciprocal-space term on the premise that the local potential already captures
short-range electrostatics. Their similar errors support the residual-design
principle, although not any particular FNO parameterization. The reported
PPPM implementation costs roughly 1.5--3 times its corresponding NEP and
scales to million-atom simulations on a single consumer GPU. The physical
source/kernel
factorization supports charge-aware observables and efficient large-scale
molecular dynamics.

Our FNO instead learns both multi-channel latent sources and the
post-deposition response operator. LES and qNEP therefore learn electrostatic
sources with a prescribed kernel. In contrast, \macefno{} learns a global
residual response operator. This flexibility could represent heterogeneous
screening or response that is
poorly compressed into a few atomic multipoles. Conversely, our model costs
parameters and mesh operations, sacrifices direct charge interpretability,
and leaves extrapolation less constrained.

\begin{table}[t]
\centering
\caption{Conceptual comparison. Entries describe the cited implementations
and this work, not every possible extension.}
\label{tab:method-comparison}
\small
\begin{tabularx}{\textwidth}{>{\bfseries}lXXX}
\toprule
 & qNEP & LES & Frozen-MACE FNO\\
\midrule
Atomic input & NEP descriptors & Local MLIP invariants &
Frozen MACE \(0e\) invariants\\
Learned source & Dynamic atomic charges & Latent charge-like variables &
Neutral multi-channel latent sources\\
Global propagation & Ewald/PPPM electrostatics & Fixed Ewald \(1/k^2\) kernel &
Learned spectral operator\\
Energy role & Explicit electrostatic term & Explicit latent-Ewald term &
Residual mesh functional\\
Electrical observables & Charges, polarization response &
Polarization and BEC demonstrated & Not identified at present\\
Geometry demonstrated & 3D periodic production MD &
3D Ewald; slab benchmark published & 2D FNO for slabs and periodic 3D FNO\\
Principal advantage & High-throughput charge-aware MD &
Compact, transferable augmentation & Flexible nonlocal response\\
Principal risk & Charge model/kernel bias & Fixed electrostatic kernel bias &
Latent gauge and extrapolation\\
\bottomrule
\end{tabularx}
\end{table}

\subsection{Quantitative context}

For water--SCAN, qNEP improves the published NEP force RMSE from 81 to
70~\meva{} in mode 2; our FNO improves its own MACE baseline from 67.57 to
52.83~\meva. The absolute values favor our model, but only a common
baseline, split, loss and parameter budget would constitute a direct
architecture comparison. qNEP additionally reports stress and electrical
response, which our current FNO calculator does not provide.

For Au$_2$/MgO, the comparison is sharper because the same DFT endpoint values
are used. MACE--LES predicts \(919.3/-56.5\)~meV for undoped/doped slabs; the
established MACE+FNO predicts \(924.5/-55.4\)~meV, compared with
\(934.8/-66.9\)~meV from DFT. Both recover the sign reversal to similar
accuracy. We therefore see the remaining scientific opportunity not as
reproducing LES with more parameters, but as constructing a tractable data set
in which spatially heterogeneous or non-electrostatic response produces a
falsifiable advantage for a learned operator.

\section{Born effective charges: opportunity and present obstacle}

For a physically defined polarization \(\mathbf P\), a Born effective charge
is
\begin{equation}
 Z^{*}_{i\alpha\beta}
 =\Omega\frac{\partial P_\alpha}{\partial R_{i\beta}}.
\end{equation}
The definition follows the density-functional perturbation-theory convention
of Ref.~\citep{gonze1997bec}.
LES and qNEP possess charge-like variables tied to a prescribed electrostatic
construction, so a polarization can be defined within the model. In our
current FNO, the channels of \(\rho_c\) admit arbitrary learned rotations and
rescalings paired with compensating changes in the operator. Neither
\(\rho_c\) nor \(\phi_c\) is calibrated to electronic charge, and the training
targets contain no polarization gauge information. A BEC obtained by simply
taking the first moment of one latent channel would therefore be
representation-dependent.

We see three plausible routes forward:
\begin{enumerate}
 \item we can retain the general residual FNO and add polarization/BEC targets that
 identify a designated physical response channel;
 \item we can split the model into an explicit electrostatic channel plus
 unconstrained residual channels;
 \item we can use BEC only as an auxiliary response diagnostic after defining a
 gauge-consistent polarization functional.
\end{enumerate}
Attributing the entire FNO residual to electrostatics would simplify the first
route, but it is a modeling assumption rather than a consequence of
energy--force training. The \(1/k^2\) audit makes that assumption more
plausible for the water model; it does not prove it.

\section{Limitations and decisive next experiments}

We identify the following limitations and corresponding next tests.

\begin{enumerate}
 \item \textbf{Continuous translation.} The residual is exactly periodic
 under grid translations but exhibits a measurable arbitrary-translation
 egg-box error. Better assignment, randomized mesh origins, analytical
 deconvolution, or efficient interlacing should be tested.
 \item \textbf{Cell heterogeneity.} A batch shares a mesh shape. Uniformly
 scaled cubes are conditioned explicitly; arbitrary anisotropic cells require
 metric-aware spectral coordinates or shape buckets.
 \item \textbf{Physical identifiability.} Latent channels are not charges.
 Electrical observables require additional structure or supervision.
 \item \textbf{Net charge.} Per-channel mean subtraction removes the
 \(k=0\) latent source and is appropriate for the neutral systems studied
 here. Charged cells and compensating-background conventions are not yet
 represented.
 \item \textbf{Stress.} Cell derivatives and virials have not been audited,
 so the ASE calculator intentionally returns only energy and forces.
 \item \textbf{Statistics.} The final models use one split and seed.
 Replicated fits are required before publication-level uncertainty claims.
 \item \textbf{Transfer.} Stable NVE/NPT trajectories,
 finite-size scaling, out-of-distribution cells and long-time observables
 remain to be demonstrated.
 \item \textbf{Deployment.} ASE inference is available. A production LAMMPS
 path needs a deployable MACE descriptor interface, particle--mesh kernels,
 a distributed FFT strategy, and validated stress.
\end{enumerate}

We see the most decisive scientific experiment as holding local environments
approximately fixed while varying a remote dielectric environment in a way
that cannot be represented by a small set of local latent multipoles. We
regard the Au$_2$/MgO sign reversal as a first instance, but need more dopant
depths, concentrations and surface terminations. Such a series would test
whether our learned response field provides information beyond a prescribed
Ewald kernel rather than merely matching it on one endpoint pair.

\section{Reproducibility and software scope}

We keep only the reusable package, numerical tests, and two benchmark
workflows in the repository. We generate training data, frozen MACE models,
caches, FNO checkpoints, Slurm logs, audits, endpoint trajectories and
compiled document artifacts outside the Git checkout. Our checkpoints store
the residual state and reconstruction metadata but reference, rather than
duplicate, the frozen MACE model.

Our test suite covers deposition, spectral energy, planar projection, the 2D
FNO slab operator, and the 3D FNO bulk operator,
cubic and planar symmetries, frozen-MACE coupling, checkpoint compatibility,
ASE evaluation, spectral diagnostics, label handling and the
adsorption-switch evaluator. We document exact commands and paths in the
benchmark READMEs. Because external model/data repositories can change, we
will additionally record upstream commit hashes, checksums, Python package
versions, Slurm scripts, random seeds and all validation indices in a
publication archive.

\section{Conclusions}

We augment a frozen local MACE potential with a conservative FNO residual
without modifying its weights. Our implementation converts invariant local
descriptors to neutral latent mesh sources, learns a geometry-aware global
response, and differentiates a scalar mesh energy. We find a substantive
reduction in held-out water--SCAN energy and force errors, and sensitivity to
a remote dopant that reverses the Au$_2$/MgO wetting preference. Our
long-wavelength audit finds an approximately \(1/k^2\) response in the cubic
water model.

Our results justify continued development, but do not yet constitute a
complete physical theory of the latent fields. LES and qNEP remain more
compact and interpretable when the missing physics is conventional
electrostatics. We expect the scientific case for FNO to be strongest where a
learned spatial response operator demonstrably captures heterogeneous
screening or other nonlocal residuals that a fixed charge--kernel
factorization misses.

\section*{Data and code availability}

We maintain the source code and reproducible benchmark launchers at
\url{https://github.com/wch3n/mace_fno}. We record results available on
1 September 2026. We do not store large third-party data, pretrained MACE
models, generated caches or learned checkpoints in the source repository; we
document their provenance and preparation commands in the benchmark
directories. We will accompany a manuscript submission with a versioned
public archive containing immutable checksums.

\bibliographystyle{unsrtnat}
\bibliography{references}
\end{document}
